\documentclass[11pt]{article}

\usepackage[T1]{fontenc}
\usepackage[utf8]{inputenc}
\usepackage{lmodern}
\usepackage[a4paper,margin=28mm]{geometry}
\usepackage{microtype}
\usepackage{amsmath,amssymb,mathtools}
\usepackage{enumitem}
\usepackage{booktabs,tabularx}
\usepackage[colorlinks=true,linkcolor=blue,citecolor=blue,urlcolor=blue]{hyperref}

\title{A Projection-First Reframing of Physics:\\
\large Descent, Effective Description, and the Limits of the Analogy}
\author{Peter Nero}
\date{July 2026\\Version 2}

\begin{document}
\maketitle

\begin{abstract}
A projection-first viewpoint asks which physical structures belong to an
effective description and which belong to a richer source from which that
description descends. This paper develops the viewpoint as a conceptual
tool, not an inevitability theorem. A many-to-one map can identify
microscopic states, but it does not by itself define lower dynamics,
probabilities, entropy, locality, an arrow of time, quantum theory, or
gravity. Those conclusions require additional structures: an intertwining
dynamics, a state or measure, observable algebras, stability estimates, and
sector-specific source maps. Modal Triplet Theory provides exact examples of
finite projectors, a shared phase line, fixed-point constructions, and
selected free field and finite Standard-Model carriers at declared tiers.
Its broader physical reconstructions remain conditional where the required
maps are incomplete. The projection-first framework is therefore most
useful as a discipline for typing levels of description and exposing missing
bridges.
\end{abstract}

\section*{Version 2 Revision Note}
\begin{description}[leftmargin=3.2cm,style=nextline]
\item[Supersedes] Version 1.0, \emph{A Projection-First Reframing of
Physics}.
\item[Reason] The original overview treated noninvertibility as sufficient
for irreversibility and sometimes presented projection-based explanations
as theorem-level necessities.
\item[Resolution] Projection, quotient dynamics, recovery, probability,
locality, and physical time are now separated. Counterexamples show why a
many-to-one map alone does not create an arrow of time.
\item[Retained result] Projection-first reasoning remains a useful way to
organize effective descriptions and to demand explicit descent and recovery
maps.
\item[Remaining boundary] Each physical application must still construct
its dynamics, state, observables, and same-source intertwiners; no universal
physics follows from projection alone.
\end{description}

\tableofcontents

\section{The basic question}

Physics often relates descriptions by forgetting detail:
\[
 \pi:X\longrightarrow Y.
\]
The source space \(X\) may contain microscopic, gauge, internal, or
preprojection data. The effective space \(Y\) contains the variables used by
an observer or lower theory.

The projection-first question is:
\[
\text{Which puzzles arise because properties of }Y
\text{ are incorrectly demanded of }X,\text{ or conversely?}
\]
This can be illuminating. Gauge redundancy, coarse-graining, reduced
density operators, effective field theories, and quotient spaces all depend
on maps between descriptive levels. Yet the notation \(\pi\) is not a
physical theory. The domain, codomain, dynamics, states, and observables must
be named.

\section{Projection and quotient dynamics}

Let \(\Phi_t:X\to X\) be upper evolution. A deterministic lower evolution
\(\psi_t:Y\to Y\) exists only if upper evolution respects projection fibers:
\[
 \pi(x)=\pi(x')
 \quad\Longrightarrow\quad
 \pi(\Phi_t x)=\pi(\Phi_t x')
\]
for every relevant \(t\). Equivalently,
\[
 \pi\Phi_t=\psi_t\pi.
\]
Without this condition, the present lower state does not determine its own
future. One may still obtain stochastic, memory-dependent, or set-valued
lower dynamics, but those require additional constructions.

This distinction is central in MTT. A finite projector can select a carrier
or sector. It does not automatically produce the effective equations acting
on that sector.

\section{Why noninvertibility is not an arrow of time}

If \(\pi\) is many-to-one, the original upper state cannot generally be
recovered from \(\pi(x)\). That is loss of distinguishability relative to
the lower description. It is not yet dynamical irreversibility.

Consider
\[
 X=Y\times F,\qquad
 \pi(y,f)=y,
\]
and an invertible lower flow \(\psi_t\). Define
\[
 \Phi_t(y,f)=(\psi_t(y),f).
\]
The projection is many-to-one whenever \(F\) has more than one point, but
the lower dynamics remains exactly reversible. Thus:
\[
\text{many-to-one description}
\not\Rightarrow
\text{irreversible lower evolution}.
\]

An arrow of time can arise when the induced lower evolution is a proper
semigroup, when recovery error grows, when entropy increases relative to a
specified state, or when a physical record process has asymmetric boundary
conditions. Each route needs its own theorem.

\subsection{Physical time and compact phase}

A periodic phase variable can function as a clock in a specified dynamical
model. It is not identical to noncompact Lorentzian time. Likewise, a branch
choice can orient a history without proving that the unobserved branch is
well behaved or physically realized.

\section{Recovery maps}

For a surjection \(\pi:X\to Y\), a section \(s:Y\to X\) satisfies
\[
 \pi s=I_Y.
\]
It is a right inverse of \(\pi\). It chooses one representative from each
fiber but does not recover an arbitrary original state, because generally
\[
 s\pi\neq I_X.
\]

Approximate recovery should be measured by a declared norm, fidelity, or
observable error. In quantum information, a recovery channel has complete
positivity and trace conditions. In geometry, a lift must preserve the
relevant connections or tensors. The word ``reconstruction'' is incomplete
until those structures are stated.

\section{Probability and quantization are additional}

A set map carries no probability by itself. Probabilities require a measure,
state, or stochastic law. Quantum probabilities additionally require an
operator algebra and positive normalized state, or an equivalent
operational structure.

Finite projectors and effects can be used inside quantum theory:
\[
 p(E|\rho)=\operatorname{tr}(\rho E).
\]
But this equation imports the trace-state pairing. It is not derived from
idempotence of \(E\). The current MTT Born-source program has exact results
on a restricted selected recorder, while the general apparatus and
objective-history problem remains open.

Similarly, a Gaussian expansion or a finite spectral decomposition does not
derive CCR/CAR quantization. The revised curved-spacetime QFT paper obtains a
selected free CAR net by composing an MTT Dirac source with independent AQFT
theorems.

\section{Locality and effective nonlocality}

An upper theory can be local in its own variables while a reduced
description is nonlocal. This happens when eliminated variables mediate
memory or when observables do not factor across the projection.

The converse warning is equally important. Calling an upper description
local does not evade Bell's theorem if the resulting model is also assumed
to satisfy measurement independence and conditional-factorization locality.
The revised Bell paper instead distinguishes upper-local equations,
microcausal observable algebras, operational no-signaling, and globally
nonfactorizing states.

For spacetime propagation, a nonlocal spatial filter need not preserve
strict support even if it leaves the hyperbolic principal symbol unchanged.
Causal claims require a retarded kernel or support estimate.

\section{What projection does organize well}

Projection-first language is especially useful for:
\begin{description}[leftmargin=3.0cm,style=nextline]
\item[Gauge theory] separating redundant representatives from physical
observables, provided constraints or BRST descent are respected;
\item[Effective theory] identifying retained variables and quantifying
discarded-sector errors;
\item[Measurement] distinguishing effects, channels, instruments, and
records;
\item[Geometry] distinguishing local frames, quotient data, and global
bundles;
\item[Fixed points] distinguishing an upper fixed state from its lower
encoding and proving any intertwiner between them;
\item[Computation] asking whether an explicit computational dynamics embeds
before invoking undecidability.
\end{description}

\section{Current MTT examples}

The present corpus contains several concrete projection structures:
\begin{itemize}
\item the \(1+2+3\) finite rank filtration;
\item the q79 finite Haar projector and complementary finite Hessian sector;
\item the shared flat phase line on the finite root-stack symbol;
\item selected finite Standard-Model carrier and profile-level transport;
\item a selected free twisted-Dirac CAR net on a globally hyperbolic
representative;
\item fixed-point and conditional-recovery theorems under explicit analytic
hypotheses.
\end{itemize}

The important open bridges include the physical nonzero-Chern HYM endpoint,
continuum geometry-to-operator naturality, selected upper action, general
Born source, nonperturbative interacting QFT, and zero-primitive
Standard-Model values.

\section{Claim ledger}

\begin{center}
\begin{tabularx}{\textwidth}{@{}
>{\raggedright\arraybackslash}p{0.32\textwidth}
>{\raggedright\arraybackslash}p{0.18\textwidth}
>{\raggedright\arraybackslash}X@{}}
\toprule
Claim & Status & Required addition\\
\midrule
Projection identifies lower-equivalent states & Exact &
Defined equivalence relation or map.\\
Compatible upper dynamics descends & Exact conditional &
Fiber-respecting intertwining equation.\\
Many-to-one projection creates an arrow of time & False in general &
Irreversible dynamics, entropy, or record theorem.\\
A section recovers the original upper state & False in general &
Section chooses a representative only.\\
Projection supplies probabilities & False in general &
Measure, state, or stochastic law.\\
Projection organizes MTT encodings & Supported &
Several exact finite examples exist.\\
All physics is inevitable from projection & Not established &
Sector-specific source maps and dynamics.\\
\bottomrule
\end{tabularx}
\end{center}

\section{Discussion}

The strongest use of a projection-first framework is methodological. It
prevents a familiar category error: transferring a property across a map
without proving that the relevant structure is preserved. It also suggests
new constructions, because a failed transfer tells us what the missing
intertwiner must carry.

For MTT, this discipline supports the proposed preprojection language. If
one upper object is to generate Dirac, Weyl, twistor, QFT, Standard-Model,
and gravitational encodings, the claim should be expressed as a family of
commuting diagrams preserving connections, actions, states, and observables.
Similarity of dimensions or terminology is not enough.

\section{Conclusion}

A projection-first reframing remains a valuable way to think about physics,
provided projection is not asked to do more than it can. It identifies
descriptive equivalence and can organize descent. Dynamics, probability,
locality, irreversibility, quantization, and gravity require additional
objects.

The current MTT program has enough exact finite structure to make this more
than a metaphor. Its next advances depend on constructing the missing
same-source intertwiners. That is a sharper and more credible program than
claiming that projection alone makes the rest inevitable.

\begin{thebibliography}{99}

\bibitem{NielsenChuang}
M.~A.~Nielsen and I.~L.~Chuang,
\emph{Quantum Computation and Quantum Information},
Cambridge University Press, 2010.

\bibitem{BreuerPetruccione}
H.-P.~Breuer and F.~Petruccione,
\emph{The Theory of Open Quantum Systems},
Oxford University Press, 2002.

\bibitem{Haag}
R.~Haag,
\emph{Local Quantum Physics},
Springer, 1996.

\bibitem{MTTFoundation}
P.~Nero,
\emph{Modal Triplet Theory: Foundations},
current revised MTT paper corpus, 2026.

\bibitem{MTTFixedPoints}
P.~Nero,
\emph{Fixed Points I--VI},
current revised MTT paper corpus, 2026.

\end{thebibliography}

\end{document}
